\chapter{总结与展望}

\section{工作总结}

本文围绕“面向 LangGraph 智能体的 AI 原生应用层故障注入机制”展开研究。随着 AI 原生应用从单模型调用演进为 Agent、工具、状态、记忆和传统微服务共同组成的复杂系统，传统资源层混沌工程和事后 AgentOps 观测都难以完整覆盖应用内部的可复现故障。TripSphere 提供了一个典型场景：旅行规划、对话修改和订单辅助分别对应线性 LangGraph workflow、ReAct 状态循环和跨 Agent A2A 委托，适合作为 AI 原生应用层可靠性验证靶场。

本文首先分析了混沌工程、AgentOps、OpenTelemetry GenAI 可观测和 LangGraph 状态图等相关工作，指出现有工具在“主动制造可复现应用层故障”方面仍存在空间。随后，本文对 TripSphere 的总体结构和三条主要链路进行了系统分析，明确链路 A 适合验证端到端规划与持久化，链路 B 适合验证多轮状态和路由鲁棒性，链路 C 适合验证跨 Agent 与跨服务可靠性。

在机制设计方面，本文提出了面向 AI 原生应用层的故障模型，将故障抽象为工具/RPC 延迟、异常、gRPC 错误、响应篡改、状态清空、路由扰动、远程 Agent 丢失和 Trace 断链等类型。围绕这些故障，本文整理了基于请求头 DSL、运行时注册表、上下文传播和 OpenTelemetry Span 属性的低侵入注入方案，使故障声明、注入命中、系统响应和实验记录能够通过 \code{experiment.id} 串联。

在靶场落地方面，TripSphere 已经形成覆盖三条链路的应用层故障注入实验闭环。本文基于 Locust artifacts 对链路 A/B/C 分别完成 baseline-low、baseline-high、fault-low 和 fault-high 四象限对比，记录了 HTTP 成功率、端到端延迟、降级信号、StateSnapshot 更新率、订单提交率等指标。结果表明，AI 原生链路中存在明显的 HTTP 可见性盲区和 fail-fast 延迟反转现象，仅依赖响应码和延迟无法判断真实业务质量。

\section{主要贡献}

本文的主要贡献可以归纳为三点。

第一，本文从系统视角分析了 AI 原生应用的可靠性问题。相比只评价模型输出，本文把 Agent 编排、工具调用、状态流转、远程 Agent 和底层微服务纳入同一系统，说明故障如何在这些边界上传播。

第二，本文设计了可复现的应用层故障注入机制。通过统一 DSL、请求级实验标识、低侵入插桩和 Span 属性记录，实验者可以对单次请求、单条 Trace 或单个工具调用施加故障，并获得可审计证据。

第三，本文给出了面向毕业论文实证部分的实验结果。四象限压测、三链路故障矩阵、业务质量指标、TraceQL 查询和链路故障案例共同构成了可复查、可复跑、可对比的实验流程。

\section{不足与局限}

本文仍有明显局限。首先，实验主要基于本地或单机环境运行，负载规模和外部 API 条件与真实生产环境仍有差距，结论更适合说明故障传播机制而非容量上限。其次，Prometheus 和 Grafana 的系统指标闭环尚不完整，容器资源、Java Actuator、spanmetrics 和业务指标还需要进一步接入。再次，LLM 输出本身具有不确定性，相同故障条件下的文本质量可能波动，后续仍需引入更稳定的结构化质量评价。最后，部分故障包含概率门，实验结果受实际命中次数影响，需要在更大样本下进一步复验。

\section{未来工作}

后续工作可以从以下方向推进。

\begin{enumerate}[label=(\arabic*)]
  \item \textbf{扩大实验样本。} 在更长时间窗口和更多并发档位下复跑三链路四象限实验，提高概率故障命中统计的稳定性。
  \item \textbf{细化故障场景。} 将当前组合故障拆分为更多单点故障实验，进一步区分 LLM、工具、状态、RPC 和 Trace 断链各自的独立影响。
  \item \textbf{完善 metrics 闭环。} 接入 cAdvisor、node-exporter、Java Actuator 和 OpenTelemetry spanmetrics，使故障命中、入口延迟、错误率和资源使用可以在 Grafana 中统一展示。
  \item \textbf{固化实验记录。} 建立 \code{experiments/} 目录，保存每次实验的 DSL、负载参数、时间窗口、PromQL、TraceQL、代表 Trace ID 和结论，提升复查与复跑能力。
  \item \textbf{探索平台化控制面。} 在当前请求头和 JSON 配置基础上，进一步设计 ConfigMap watcher、CRD 或与 Chaos Mesh 的集成，使应用层故障注入可以进入更标准的混沌工程工作流。
\end{enumerate}

总体而言，本文完成了 TripSphere 作为 AI 原生应用层故障注入靶场的架构分析、机制设计、工程落地和首轮实验验证。后续随着实验样本扩大、指标闭环完善和平台化控制面接入，该工作可以进一步发展为面向 LangGraph 智能体可靠性验证的可复现 Benchmark。
